\section{Future Work}

Several directions remain unexplored.

\paragraph{Prompt format optimization.} We use a fixed prompt template throughout and do not attempt to optimize it for ICL elicitation.
Prompt engineering could improve the metric's discriminative power.
Dimension-specific framings (e.g., ``Here are several stories in different genres'') are one such variant: priming $\theta$ to expect variation along a dimension removes the small portion of the surprise attributable to the mere existence of that variation, but it does not isolate per-dimension diversity, since the rest of the cross-response surprise reduction remains.

\paragraph{Ensembling base models.} A single base model $\theta$ may produce non-monotone $a_k$ curves when pushed out of distribution by long contexts.
Ensembling multiple base models at the token level (averaging softmax probabilities) could stabilize the estimates while preserving the autoregressive structure.
Our codebase supports token-level ensembling, but we have not yet validated it experimentally.

\paragraph{Total-bits variants of $D$.}
We report $D_{Ca_n} = C \times a_n$ in bits/byte for length control, normalising by bytes rather than tokens so that scores remain tokenizer-agnostic.
The implementation also exposes $a_n$ in total bits, which is likewise tokenizer-agnostic but not length-controlled.
A hybrid scalar $C \times a_n^{\text{bits}}$ rewards length (longer responses have more positions at which to be surprising), which some practitioners may want; characterising when each variant is the appropriate measurement target is left to future work.

